% ASSERT_CONVENTION: natural_units=dimensionless, metric_signature=mostly_minus,
%   jordan_product=(1/2)(ab+ba), octonion_basis=fano_e1e2=e4,
%   complex_structure=u_equals_e7, peirce_decomposition=under_E11,
%   det3_normalization=d(X,X,X)=6*det_3(X),
%   peirce_basis_ordering=I0_V1_I1to16_Vhalf_I17to26_V0,
%   spacetime_V0_indices={17,18,19,26}, internal_V0_indices={20,...,25},
%   C_IJK=(1/6)*d_IJK

\documentclass[11pt]{article}
\usepackage{amsmath,amssymb,amsthm}
\newtheorem{proposition}{Proposition}
\newtheorem{theorem}{Theorem}
\newtheorem{lemma}{Lemma}
\newtheorem{remark}{Remark}
\newcommand{\Tr}{\mathrm{Tr}}
\newcommand{\h}{\mathfrak{h}}
\begin{document}

\section{Stress-Energy Identification and Weinberg Theorem Application}

\subsection{Setup and Notation}

From Phase 49: the coupling tensor $C_{IJK} = \frac{1}{6}\, d_{IJK}$ decomposes
into two nonzero Peirce blocks (Phase 47):
\begin{itemize}
\item $(V_1, V_0, V_0)$: 10 entries (gravitational self-coupling via $\det_2$)
\item $(V_{1/2}, V_{1/2}, V_0)$: 96 entries (matter-gravity coupling)
\end{itemize}
The 96 matter-gravity entries split further by the $V_0$ index (Phase 49):
\begin{itemize}
\item Spacetime ($a \in \{17,18,19,26\}$): 48 unique triples
\item Internal ($a \in \{20,\ldots,25\}$): 48 unique triples
\end{itemize}
We analyze the spacetime couplings $C_{i,j,a}$ with $i,j \in V_{1/2}$
(Peirce indices 1--16, the 16 matter fields) and $a$ in the spacetime
$V_0$ set, restricting to the metric perturbation $h_{ab}$ on
$\h_2(\mathbb{C}_u) = \mathbb{R}^{3,1}$.

\subsection{Stress-Energy Identification}

\begin{proposition}[Stress-energy coupling structure]\label{prop:stress-energy}
The spacetime coupling $C_{i,j,a}$, restricted to $a \in \{17,18,19,26\}$
(spacetime $V_0$), has the structure of a stress-energy coupling:
\begin{enumerate}
\item \textbf{Symmetry:} $C_{i,j,a} = C_{j,i,a}$ for all
  $i,j \in V_{1/2}$ and all spacetime~$a$
  (verified exactly: $\max |C_{i,j,a} - C_{j,i,a}| = 0$, 480 pairs checked).
\item \textbf{Universality:} All 16 matter fields couple to all 4
  spacetime directions. Each matter field $\phi^i$ ($i = 1,\ldots,16$)
  has nonzero couplings $C_{i,j,a}$ for \emph{every} spacetime
  direction $a \in \{17,18,19,26\}$.
  Per-index unique triple counts: $17{:}16$, $18{:}16$, $19{:}8$, $26{:}8$
  (total 48, matching Phase 49).
\item \textbf{Bilinear in matter:} The coupling $C_{i,j,a}\,\phi^i\,\phi^j\,h_{ab}$
  is bilinear in the matter fields $\phi^i$ (from $V_{1/2}$) and linear
  in the metric perturbation $h_{ab}$ (from spacetime $V_0$).
\end{enumerate}
\end{proposition}

\begin{proof}
\textbf{Symmetry} follows from the total symmetry $C_{IJK} = C_{(IJK)}$
of the coupling tensor (Phase 47: $d_{IJK}$ is fully symmetric,
verified over 50 random triples with max error 0).
Verified explicitly for all 480 ordered pairs
$(i,j)$ with $i < j$ across 4 spacetime directions:
$\max |C_{i,j,a} - C_{j,i,a}| = 0$ (exact to machine precision).

\textbf{Universality} was established in Phase 49
(decompose\_couplings\_49: all 4 spacetime $V_0$ indices present,
all 16 matter indices present) and re-verified here. Each matter
field couples to all 4 spacetime directions (4 nonzero couplings per
matter field). The coupling structure is:
\begin{itemize}
\item $a = 17$ (timelike $b_0$): $C_{i,i,17} = -1/6$ for all $i = 1,\ldots,16$
  (diagonal, universal, equal strength).
\item $a = 18$ (spatial $b_1$): $C_{i,i,18} = -1/6$ for $i = 1,\ldots,8$;
  $C_{i,i,18} = +1/6$ for $i = 9,\ldots,16$ (diagonal, sign split).
\item $a = 19$ (spatial $b_2 = e_0$): 8 off-diagonal entries
  $C_{i,j,19}$ with $|C| = 1/3$, pairing matter fields $\{1,2,\ldots,8\}$
  with $\{9,10,\ldots,16\}$.
\item $a = 26$ (spatial $b_9 = e_7$): 8 off-diagonal entries
  $C_{i,j,26}$ with $|C| = 1/3$, pairing the same sets with
  different pairings (related to the complex structure $u = e_7$).
\end{itemize}

\textbf{Bilinearity} is immediate from the form of the interaction:
$\mathcal{L}_{\mathrm{int}} \supset C_{i,j,a}\,\phi^i\,\phi^j\,h^a$,
which is quadratic in $\phi$ and linear in $h$.
\end{proof}

\begin{remark}[Trace coupling]\label{rem:trace}
The trace coupling $T_{ij} = \sum_{a \in \mathrm{spacetime}} G^{aa}\,C_{i,j,a}$,
where $G^{aa}$ is the inverse $\det_2$ Gram on the Peirce basis
($G^{-1} = \mathrm{diag}(+4, -4, -1, -1)$ for $\{17,18,19,26\}$),
is nonzero with Frobenius norm $\|T\|_F \approx 4.22$.
This matrix has:
\begin{itemize}
\item Diagonal entries $T_{ii} = 0$ for $i = 1,\ldots,8$ (cancellation
  between the equal-magnitude but opposite-sign contributions from
  $a = 17$ and $a = 18$, with $G^{-1}_{17} = +4$ and $G^{-1}_{18} = -4$).
\item Diagonal entries $T_{ii} = -4/3$ for $i = 9,\ldots,16$
  (the $a = 18$ contribution changes sign, so the two diagonal
  couplings reinforce rather than cancel).
\item 16 off-diagonal entries from $a = 19$ and $a = 26$ with
  $|T_{ij}| = 1/3$.
\end{itemize}
The nonzero trace coupling confirms that matter couples to both
the spin-2 (traceless) and spin-0 (trace) parts of $h_{ab}$,
as expected for a stress-energy coupling. In the massless spin-2
theory, the spin-0 mode decouples from the S-matrix
(van Dam--Veltman--Zakharov), but the algebraic coupling is present.
\end{remark}

\begin{remark}[Non-derivative identification]\label{rem:nonderivative}
The algebraic coupling $C_{i,j,a}\,\phi^i\,\phi^j\,h^a$ provides the
\textbf{non-derivative part} of the stress-energy coupling.
The full stress-energy tensor $T_{ab}$ for scalar fields includes
derivative terms:
\[
  T_{ab} = \partial_a \phi^i \partial_b \phi_i
    - \tfrac{1}{2}\,\eta_{ab}\,\partial_c \phi^i \partial^c \phi_i
    + \text{potential terms}.
\]
The derivative couplings $({\partial_a \phi}\,{\partial_b \phi}\,h^{ab})$
arise from \textbf{minimal coupling}: replacing $\eta_{ab} \to g_{ab} = \eta_{ab} + h_{ab}$
in the kinetic term $g^{ab}\partial_a\phi\,\partial_b\phi$.
This is the standard prescription once the metric perturbation is identified.

What the $\h_3(\mathbb{O})$ algebraic structure provides is:
\begin{itemize}
\item The non-derivative coupling $C_{i,j,a}\,\phi^i\phi^j\,h^a$ (from $d_{IJK}$);
\item The identification of $h_{ab}$ as a metric perturbation
  with $\det_2$ kinetic structure (from Propositions 1--2 in Plan 01);
\item The universality of the coupling (all 16 matter fields, all 4
  spacetime directions).
\end{itemize}
The derivative terms arise from minimal coupling (covariantization),
which is the unique consistent coupling prescription for a metric field.
\end{remark}

\subsection{Weinberg Theorem Application}

\begin{theorem}[GR from $\h_3(\mathbb{O})$ via Weinberg]\label{thm:weinberg}
The algebraic structure of $\h_3(\mathbb{O})$ --- specifically, the Peirce
decomposition $27 = 1 + 16 + 10$ with its cubic form $\det_3$, the
bilinear $\det_2$ on $V_0$, and the coupling tensor $C_{IJK}$ ---
satisfies all four hypotheses of Weinberg's theorem (1964). Therefore,
the Einstein--Hilbert term $-R/2$ is forced at low energies.
\end{theorem}

\begin{proof}
We verify each of the four Weinberg hypotheses with explicit algebraic
sources, then apply the theorem.

\medskip
\noindent\textbf{H1. Lorentz invariance.}
\emph{Source:} Phase 48.
The compact group $\mathrm{Spin}(9) \subset \mathrm{Aut}(\h_3(\mathbb{O}))$
stabilizes $V_0$. The stabilizer subalgebra decomposes as
$\mathfrak{so}(3) \oplus \mathfrak{so}(6)$ (dim 18), where $\mathfrak{so}(3)$
acts on the 4 spacetime $V_0$ directions as the rotation subalgebra of
$\mathfrak{so}(3,1)$. The full Lorentz algebra is recovered via
complexification: $\mathfrak{so}(3,\mathbb{C}) \cong \mathfrak{sl}(2,\mathbb{C})$,
which is the complexified Lorentz algebra. The projection $\pi_u$ is
equivariant under the full 18-dimensional stabilizer (max error
$5.15 \times 10^{-17}$, Phase 48).

\emph{Non-circularity:} The Lorentz structure arises from the
$F_4/\mathrm{Spin}(9)$ structure of $\h_3(\mathbb{O})$, not from
assuming GR or the Einstein field equations.

\emph{Caveat:} The compact $\mathrm{Spin}(9)$ contains only the
rotation subgroup $\mathrm{SO}(3)$, not the full non-compact
$\mathrm{SO}(3,1)$. Boosts are absent from the compact group.
The full Lorentz group is recovered by analytic continuation
($\mathfrak{so}(3,\mathbb{C}) = \mathfrak{sl}(2,\mathbb{C})$
contains boosts after complexification). This is the standard
relation between compact and non-compact real forms of the same
complexified Lie algebra (see Phase 48 for details).

\medskip
\noindent\textbf{H2. Spin-2.}
\emph{Source:} Plan 01, Proposition~1.
The symmetric bilinear perturbation $h_{ab}$ of $\det_2$ on
$\h_2(\mathbb{C}_u) = \mathbb{R}^{3,1}$ decomposes under
$\mathrm{SO}(3,1)$ as
$10 = 9\,(\text{spin-2, traceless symmetric}) + 1\,(\text{spin-0, trace})$.
The traceless projector has rank 9 (verified: idempotent error $< 10^{-15}$,
orthogonal and complete with trace projector of rank 1).

\emph{Non-circularity:} The spin-2 identification comes from
representation theory of the $\det_2$ perturbation on $\h_2(\mathbb{C}_u)$.
It does not use the Einstein field equations or assume $-R/2$ in the action.

\medskip
\noindent\textbf{H3. Massless.}
\emph{Source:} Plan 01, Proposition~2.
The $\det_3$ expansion around the rank-1 idempotent $E_{11}$ gives
\[
  \det_3(E + \varepsilon\delta) = \varepsilon^2\,\det_2(\delta)
    + \varepsilon^3\,\det_3(\delta),
\]
where $\det_3(\delta) = 0$ for all $\delta \in V_0$ (structural zero).
The quadratic coefficient $M_{ab} = \det_2$ is the metric norm-squared
(kinetic-type), not the Fierz--Pauli mass
$m^2(h_{ab}h^{ab} - h^2)$. No mass term exists.
Verified by analytical computation (polarized sharp map) and numerical
finite differences (agreement $< 2 \times 10^{-16}$), with independent
confirmation via $F_4$-invariant decomposition.

\emph{Non-circularity:} Masslessness follows from the algebraic
structure of $\det_3$ (specifically $E^\# = 0$ for the rank-1
idempotent $E_{11}$). It does not use $-R/2$ or the Einstein equations.
The key identity $\Tr(\delta^\# \circ E) = \det_2(\delta)$ is a
Jordan algebra identity, not a gravitational equation of motion.

\medskip
\noindent\textbf{H4. Universal coupling to stress-energy.}
\emph{Source:} This plan, Proposition~3.
The spacetime coupling $C_{i,j,a}$ is symmetric ($C_{i,j,a} = C_{j,i,a}$,
exact), universal (all 16 matter fields couple to all 4 spacetime
directions), and bilinear in matter fields contracted with the metric
perturbation. The trace coupling $T_{ij} = G^{aa} C_{i,j,a}$ is
nonzero ($\|T\|_F \approx 4.22$), confirming coupling to both spin-2
and spin-0 parts of $h_{ab}$.

The coupling $C_{i,j,a}\,\phi^i\phi^j\,h^a$ provides the non-derivative
part of the stress-energy coupling (Remark~\ref{rem:nonderivative}).
The derivative part arises from minimal coupling (covariantization),
which is the unique consistent prescription for a massless spin-2
metric field.

\emph{Non-circularity:} The universality of the coupling comes from
the $C_{IJK} = \tfrac{1}{6}\,d_{IJK}$ decomposition (Phase 49),
which traces to the polarization of $\det_3$ on $\h_3(\mathbb{O})$.
No assumption of $-R/2$ or the Einstein field equations is used.

\medskip
\noindent\textbf{Application of Weinberg's theorem.}
All four hypotheses are verified from the algebraic structure of
$\h_3(\mathbb{O})$. By Weinberg's theorem (Weinberg 1964,
\emph{Phys.\ Rev.}\ \textbf{135}, B1049):

\begin{quote}
In any Lorentz-invariant quantum field theory containing a massless
spin-2 field with universal coupling to the stress-energy tensor,
the low-energy effective theory must reproduce general relativity.
\end{quote}

Therefore, the Einstein--Hilbert term $-R/2$ is \textbf{forced}
at low energies by the algebraic structure of $\h_3(\mathbb{O})$.
Combined with the prepotential $F(X) = \det(X)$ from Phase 49
(which determines the scalar manifold, gauge kinetic matrix, and
all matter-gravity couplings), the complete bosonic Lagrangian
(Eq.~49.6) is determined:
\begin{equation}\label{eq:50.9}
  e^{-1}\,\mathcal{L}_{\mathrm{bos}}
    = \underbrace{-\frac{R}{2}}_{\text{Weinberg}}
    + \underbrace{g_{i\bar{j}}\,\partial_\mu z^i \partial^\mu \bar{z}^{\bar{j}}
    + \mathrm{Im}(\mathcal{N}_{IJ})\, F^I_{\mu\nu} F^{J\mu\nu}
    + \mathrm{Re}(\mathcal{N}_{IJ})\, F^I_{\mu\nu} {*F}^{J\mu\nu}}_{\text{from }\det(X)}.
\end{equation}
\end{proof}

\begin{remark}[Non-circularity summary]\label{rem:noncircular}
None of the four Weinberg inputs uses the Einstein field equations,
the Einstein--Hilbert action, or the assumption that $-R/2$ appears
in the Lagrangian:
\begin{enumerate}
\item[H1:] $\mathrm{Spin}(9) \subset \mathrm{Aut}(\h_3(\mathbb{O}))$
  $\to$ stabilizer of $V_0$ $\to$ $\mathfrak{so}(3) \oplus \mathfrak{so}(6)$
  $\to$ $\mathfrak{so}(3,1)$ via complexification.
  \emph{Source: $F_4/\mathrm{Spin}(9)$ structure, not GR.}
\item[H2:] $\det_2$ Gram $= \mathrm{diag}(+1,-1,-1,-1)$ $\to$
  symmetric perturbation $\to$ $\mathrm{SO}(3,1)$ irrep.
  \emph{Source: Jordan algebra $\det_2$, not GR.}
\item[H3:] $\det_3(E + \varepsilon\delta) = \varepsilon^2\det_2(\delta)$
  $\to$ $E^\# = 0$ $\to$ no quadratic potential.
  \emph{Source: Jordan algebra $\det_3$ rank structure, not GR.}
\item[H4:] $C_{IJK} = \tfrac{1}{6}\,d_{IJK}$ $\to$
  $(V_{1/2},V_{1/2},V_0)$ block $\to$ symmetric, universal.
  \emph{Source: Jordan algebra cubic form polarization, not GR.}
\end{enumerate}
The conclusion $-R/2$ is \textbf{derived} from the algebraic
structure of $\h_3(\mathbb{O})$ via Weinberg's theorem, not assumed.
\end{remark}

\begin{remark}[Scope and caveats]\label{rem:scope}
Weinberg's theorem is a low-energy/perturbative result. It establishes
that $-R/2$ is the unique \emph{low-energy effective} gravitational
action consistent with the four hypotheses. It does not constrain:
\begin{itemize}
\item UV completion or non-perturbative gravitational effects;
\item Higher-derivative corrections (e.g., $R^2$, $R_{\mu\nu}R^{\mu\nu}$)
  that are suppressed at low energies;
\item The cosmological constant ($\Lambda = 0$ for ungauged MESGT,
  from Phase 49, Proposition 6).
\end{itemize}
Additionally, the stress-energy identification is for the
non-derivative coupling; derivative terms arise from minimal
coupling (Remark~\ref{rem:nonderivative}).
\end{remark}

\end{document}
